\documentclass{report}
\usepackage{amsmath,amssymb}
\setlength{\parindent}{0mm}
\setlength{\parskip}{1em}
\begin{document}
\begin{center}
\rule{6in}{1pt} \
{\large Michael Griebel \\
{\bf Greedy and Randomized Versions of the Multiplicative Schwarz Method}}

Institut f\"ur Numerische Simulation \\ Universit\"at Bonn \\ Wegelerstr 6 \\ D-53115 Bonn \\ and \\ Fraunhofer Institute for Algorithms and Scientific Computing SCAI \\ Schloss Birlinghoven \\ D-53754 Sankt Augustin
\\
{\tt griebel@ins.uni-bonn.de}\end{center}

We consider sequential, i.e., Gauss-Seidel type, subspace correction
methods for the iterative solution of symmetric positive definite
variational problems, where the order of subspace correction steps is
not deterministically fixed as in standard multiplicative Schwarz
methods. Here, we greedily choose the subspace with the largest (or at
least a relatively large) residual norm for the next update step, which
is also known as the Gauss-Southwell method. We prove exponential
convergence in the energy norm, with a reduction factor per iteration
step directly related to the spectral properties, e.g., the condition
number, of the underlying space splitting. To avoid the additional
computational cost associated with the greedy pick, we alternatively
consider choosing the next subspace randomly, and show similar estimates
for the expected error reduction. We give some numerical examples, in
particular applications to a Toeplitz system and to multilevel
discretizations of an elliptic boundary value problem, which illustrate
the theoretical estimates.

This is joint work with Peter Oswald (JU Bremen).


\end{document}
